% !TEX program = xelatex
% Compile with:  xelatex g3_model_theory.tex   (run twice for the ToC / refs)
% Requires XeLaTeX + a Traditional-Chinese system font (Microsoft JhengHei).
\documentclass[11pt]{article}

\usepackage{fontspec}
\usepackage{xeCJK}
\setCJKmainfont{Microsoft JhengHei}
% Latin main/mono fonts left at the fontspec defaults (Latin Modern) for portability.

\usepackage[a4paper,margin=1in]{geometry}
\usepackage{amsmath,amssymb}
\usepackage{booktabs}
\usepackage{array}
\usepackage{longtable}
\usepackage{xcolor}
\usepackage{enumitem}
\usepackage{titlesec}
\usepackage[hidelinks,colorlinks=true,linkcolor=blue!55!black,urlcolor=teal!60!black,citecolor=blue!55!black]{hyperref}

\definecolor{g2col}{HTML}{1f6feb}
\definecolor{g3col}{HTML}{d9480f}
\definecolor{codebg}{HTML}{f4f4f4}

\newcommand{\rv}{\mathbf{r}}
\newcommand{\Fv}{\mathbf{F}}
\newcommand{\dt}{\Delta t}
\newcommand{\umt}{\mu\text{m}}            % micrometre, rendered in math mode
\newcommand{\um}{\,\umt}
\newcommand{\nN}{\,\text{nN}}
\newcommand{\GTwo}{\textcolor{g2col}{\textbf{G2}}}
\newcommand{\GThree}{\textcolor{g3col}{\textbf{G3}}}
\newcommand{\code}[1]{\texttt{\small #1}}

\titleformat{\section}{\Large\bfseries\color{black!85}}{\thesection}{0.6em}{}
\titleformat{\subsection}{\large\bfseries\color{black!75}}{\thesubsection}{0.6em}{}

\title{\vspace{-1.2cm}\textbf{The Generation-3 Spheroid-Guidance Model}\\[2pt]
\large Theory, Equations, and What Changed From Generation~2\\[6pt]
\normalsize\textnormal{Motor--clutch collagen-remodelling project \;/\; \code{generations/g3\_spheroid\_guidance}}}
\author{}
\date{2026-08-19}

\begin{document}
\maketitle
\vspace{-1.0cm}

% ======================================================================
\begin{abstract}
\noindent\textbf{中文摘要。}\;
Generation~3（\GThree，\code{generations/g3\_spheroid\_guidance}）是一個\textbf{二維機制示範}模型，
用來說明一顆癌細胞\textbf{spheroid（球體）}如何在\emph{沒有任何預先指定方向}的情況下，
自發打破對稱、伸出 protrusion（偽足）跨越一圈「無纖維間隙」去抓住 collagen（膠原）纖維、
用 motor--clutch（馬達--離合器）張力把周圍纖維網路往內拉並重排，最後往它真正抓得住的那一側遷移。
\GThree\ \textbf{原封不動地重用了} Gloria 已驗證的 Generation~2（\GTwo）膠原力學引擎
（\code{generations/g2\_corrected/common/model.py}）：珠--彈簧網路、過阻尼動力學、
軸向彈簧、離散彎曲、可自由轉動的永久 crosslink（交聯）、邊界錨定與細胞表面軟排斥全部不變。
\GThree\ 只加了四樣東西：（1）帶「無纖維間隙」的 spheroid 幾何；（2）顯式 protrusion 與
tip-first（尖端先接觸）綁定；（3）\textbf{emergent polarity（湧現極性）}——一個質量守恆的
wave-pinning（波釘）活性場加上 FAK/Rac1 式的黏著回饋，\textbf{取代了}\GTwo\ V3 寫死的
\code{polarity\_probability = 0.65}；（4）由反作用力驅動的\emph{二維}遷移。
本文用中英對照逐一寫出所有方程式、說明其運作，並明確標出哪些與\GTwo\ 相同、哪些被改動。

\medskip
\noindent\textbf{English.}\;
\GThree\ is a two-dimensional \emph{mechanism demonstration}: it asks whether a cancer-cell
\emph{spheroid} can, with no prescribed direction, break symmetry, extend protrusions across an
initially fibre-free gap to grip collagen, remodel the surrounding network by motor--clutch
traction, and migrate toward the side where it actually grips. \GThree\ \emph{reuses the frozen,
validated \GTwo\ collagen engine unchanged} and adds only four ingredients: a spheroid with a
fibre-free gap, explicit tip-first protrusions, an \emph{emergent} wave-pinning polarity field
with FAK/Rac1 adhesion feedback (which \emph{replaces} the hardcoded
\code{polarity\_probability = 0.65} of \GTwo\ V3), and reaction-driven 2-D migration. This report
states every equation, explains how it works, and marks precisely what is shared with \GTwo\ and
what changed. \emph{This is not a calibrated 3-D tumour-migration prediction.}
\end{abstract}

\tableofcontents
\vspace{4mm}

% ======================================================================
\section{Positioning and lineage / 定位與傳承}

The project is organised as a sequence of ``generations''. Each generation freezes the previous
one and isolates a single new mechanism.

\begin{itemize}[leftmargin=1.4em,itemsep=2pt]
  \item \GTwo\ \emph{(corrected)} --- a crosslinked, boundary-anchored bead--spring collagen
        network with overdamped dynamics, validated at $180\um$ / 99 fibres / $5\nN$. Its
        sub-versions are \textbf{V2} (elastic crosslink transmission, fixed cell),
        \textbf{V3} (two-sided motor clutches $+$ rigid-cell translation, $\sim\!0.28\um$/min
        migration), and \textbf{V4} (contact plasticity, a negative baseline).
  \item \GThree\ \emph{(spheroid guidance)} --- the subject of this report. It reuses the \GTwo\
        \emph{engine} verbatim and adds a spheroid, protrusions, emergent polarity, and
        2-D reaction-driven motion.
\end{itemize}

\noindent\textbf{Why a rebuild.} An earlier \GThree\ (\code{src/g3/}, now archived under
\code{legacy/g3\_v1\_superseded/}) built a brand-new spatial-clutch engine from scratch. It was
mis-scaled, ran for $15$--$60\,\text{s}$ of simulated time --- far too short for the
tens-of-minutes-to-hours over which fibre realignment and migration actually develop --- and
produced chaotic animations that did not match \GTwo. Its guidance test even \emph{failed}
(aligned nematic order $0.0142 <$ isotropic $0.0160$). The current \GThree\ discards that engine and
instead stands on the already-validated \GTwo\ mechanics.

\medskip
\noindent\textbf{Units.} Throughout, lengths are micrometres ($\umt$), forces are
nanonewtons ($\text{nN}$), and time is seconds ($\text{s}$) --- identical to \GTwo, so the two
models share one coordinate and force language.

% ======================================================================
\section{The shared Generation-2 engine (unchanged) / 共用的 G2 引擎（不變）}
\label{sec:engine}

Everything in this section is \emph{inherited from \GTwo\ without modification}. \GThree's
\code{SpheroidConfig} subclasses \GTwo's \code{CollagenConfig}, and \GThree's integrator
\code{fast\_advance} reproduces \GTwo's \code{Network.advance} bit-for-bit (it only swaps
\code{np.add.at} for \code{np.bincount} scatter-adds, $\sim\!2.4\times$ faster). The collagen
physics below is therefore common ground.

\subsection{Overdamped bead force balance / 過阻尼力平衡}

Collagen is a set of bead chains. Every mobile bead $i$ obeys the professor's overdamped
(inertia-free) force balance
\begin{equation}
  \zeta\,\dot{\rv}_i \;=\;
  \Fv_i^{\text{stretch}} + \Fv_i^{\text{bend}} + \Fv_i^{\text{crosslink}}
  + \Fv_i^{\text{repulsion}} + \Fv_i^{\text{active}},
  \label{eq:balance}
\end{equation}
where $\zeta$ is the bead drag (\code{bead\_drag}). The viscous term $\zeta\dot\rv_i$ dissipates
motion into the surrounding fluid/porous matrix; the elastic bonds store mechanical energy.
There is deliberately \textbf{no} standard-linear-solid (SLS) element and \textbf{no} plasticity.
The equation is integrated with an explicit forward-Euler step,
\begin{equation}
  \rv_i(t+\dt) = \rv_i(t) + \dt\,\frac{\Fv_i^{\text{total}}}{\zeta}
  \quad (i \notin \text{fixed}),
  \qquad \rv_i \equiv \rv_i^{0}\ \ (i \in \text{fixed}).
  \label{eq:euler}
\end{equation}
Boundary beads (those with $\max(|x_i|,|y_i|)\ge \tfrac{L}{2}-w_b$, where $L$ is the domain size
and $w_b$ the boundary band width) are \emph{fixed}: their positions are pinned to $\rv_i^0$ so
the network is anchored to a rigid outer frame. Stability requires $\dt < 2\zeta/k_t$; \GThree\ uses
$\dt = 0.05\,\text{s}$ against $2\zeta/k_t \approx 0.32\,\text{s}$.

\subsection{Axial stretch / 軸向拉伸}

Each bond (edge) joins consecutive beads $i,j$ with initial rest length $\ell_0=\|\rv_j^0-\rv_i^0\|$.
Writing $\mathbf d=\rv_j-\rv_i$, $\ell=\|\mathbf d\|$, extension $e=\ell-\ell_0$, the bond is a
Hookean spring whose stiffness switches between tension and compression:
\begin{equation}
  k =
  \begin{cases}
    k_t = EA/\ell_0, & e \ge 0 \quad(\text{tension}),\\[2pt]
    k_c = \rho\,k_t,  & e < 0 \quad(\text{compression / microbuckling}),
  \end{cases}
  \qquad
  \Fv_i^{\text{stretch}} = +k\,e\,\hat{\mathbf d},\quad
  \Fv_j^{\text{stretch}} = -k\,e\,\hat{\mathbf d}.
  \label{eq:stretch}
\end{equation}
Here $EA$ is the axial rigidity (\code{axial\_rigidity}), $EA = E\,A$ with fibre cross-section
$A=\pi(D/2)^2$, diameter $D$, and modulus $E$. The compression ratio $\rho=0.10$
(\code{compression\_ratio}) softens compressed bonds, representing fibre microbuckling
(Notbohm 2015). The stored energy is $U_{\text{stretch}}=\tfrac12\sum k\,e^2$.

\subsection{Discrete bending / 離散彎曲}

For every interior triplet $(a,b,c)$ of consecutive beads on a fibre, the discrete curvature is
$\boldsymbol\kappa = \rv_a - 2\rv_b + \rv_c$ with a stress-free reference $\boldsymbol\kappa_0$
taken from the initial geometry. Bending penalises the deviation
$\Delta\boldsymbol\kappa=\boldsymbol\kappa-\boldsymbol\kappa_0$:
\begin{equation}
  \Fv_a^{\text{bend}} = -k_b\,\Delta\boldsymbol\kappa,\quad
  \Fv_b^{\text{bend}} = +2k_b\,\Delta\boldsymbol\kappa,\quad
  \Fv_c^{\text{bend}} = -k_b\,\Delta\boldsymbol\kappa,
  \qquad
  k_b = \frac{EI}{\bar\ell^{\,3}},
  \label{eq:bend}
\end{equation}
where $EI$ is the bending rigidity (\code{bending\_rigidity}), $EI = E\,I$ with second moment
$I=\pi D^4/64$, and $\bar\ell$ is the mean bond length. These forces are the negative gradient of
$U_{\text{bend}}=\tfrac12 k_b\sum\|\Delta\boldsymbol\kappa\|^2$.

\subsection{Freely-hinged permanent crosslinks / 可自由轉動的永久交聯}

Where two \emph{different} fibres cross, a permanent hinged link joins the two coincident
material points. A material point on edge $(a_1,a_2)$ at parameter $\alpha_a$ is
$\mathbf p_a=(1-\alpha_a)\rv_{a_1}+\alpha_a\rv_{a_2}$ (similarly $\mathbf p_b$). The link stores
\begin{equation}
  U_x = \tfrac12 k_x\,\bigl\|(\mathbf p_b-\mathbf p_a)-\mathbf d_{x,0}\bigr\|^2,
  \qquad
  \mathbf f = k_x\bigl[(\mathbf p_b-\mathbf p_a)-\mathbf d_{x,0}\bigr],
  \label{eq:xlink}
\end{equation}
with $\mathbf d_{x,0}=\mathbf 0$ for a permanent link (the two points are initially coincident)
and stiffness $k_x$ (\code{crosslink\_stiffness}). The force $\mathbf f$ is distributed back onto
the four host beads by the interpolation weights:
$+(1-\alpha_a)\mathbf f$ on $a_1$, $+\alpha_a\mathbf f$ on $a_2$,
$-(1-\alpha_b)\mathbf f$ on $b_1$, $-\alpha_b\mathbf f$ on $b_2$. The link is \emph{freely hinged}
(it resists separation of the two points, not the angle between fibres) and \emph{never breaks} in
\GTwo/\GThree\ --- there is no crosslink rupture manufacturing motion.

\subsection{Soft cell-surface repulsion / 細胞表面軟排斥}

The cell is a rigid circle of radius $R$ centred at $\rv_c$. Beads that penetrate its surface
(plus a small clearance $c$) feel a one-sided linear penalty that pushes them radially out:
\begin{equation}
  \delta_i = \max\!\bigl(0,\; R + c - \|\rv_i-\rv_c\|\bigr),
  \qquad
  \Fv_i^{\text{repulsion}} = k_{\text{rep}}\,\delta_i\,
  \frac{\rv_i-\rv_c}{\|\rv_i-\rv_c\|},
  \label{eq:repulsion}
\end{equation}
with $k_{\text{rep}}$ = \code{repulsion\_stiffness} and $c$ = \code{cell\_clearance}. This is the
repulsive branch of a linear-elastic contact; it adds no adhesion and no friction.

\subsection{Network generation and the percolation gate / 網路生成與滲流閘}
\label{sec:gen}

Finite curved fibres of length $20$--$110\um$ are laid down (a sine-perturbed segment, function
\code{\_curved\_segment}, reused from \GTwo). Interior fibres are seeded at random angles; a fraction
are ``boundary-seeded'' so they reach the anchored frame. Permanent crosslinks are then placed at
all pairwise fibre intersections (\code{build\_crosslinks}). A generated network is \textbf{accepted
only if} it percolates: at least a required fraction of fibres are connected through the crosslink
graph to the fixed boundary (\code{connectivity\_report}). \GThree\ relaxes the gate slightly
(required connected fraction $0.80$ vs $0.85$) and, crucially, changes the geometry --- see
\S\ref{sec:gap}.

% ======================================================================
\section{The shared motor--clutch law / 共用的馬達--離合器定律}
\label{sec:clutch}

The stochastic motor--clutch law is also inherited from \GTwo\ V3 (Bell 1978; Bangasser \& Odde
2013). \GThree\ applies \emph{the identical law}, but per \emph{protrusion} rather than to two fixed
left/right ensembles. A bundle has $C$ parallel Hookean clutches; clutch $c$ has extension
$x_c$ and force $F_c=k_c x_c$ (stiffness $k_c$ = \code{clutch\_stiffness}).

\paragraph{Bond kinetics (Bell slip-bond).}
An unbound clutch binds within $\dt$ with probability $1-e^{-k_{\text{on}}\dt}$. A bound clutch's
off-rate accelerates with its own load, and it unbinds with probability
$1-e^{-k_{\text{off}}\dt}$:
\begin{equation}
  k_{\text{off}} = k_{\text{off}}^{0}\,\exp\!\Bigl(\tfrac{F_c}{F_b}\Bigr),
  \qquad
  P_{\text{unbind}} = 1-e^{-k_{\text{off}}\dt},
  \qquad
  P_{\text{bind}} = 1-e^{-k_{\text{on}}\dt},
  \label{eq:bell}
\end{equation}
with unloaded off-rate $k_{\text{off}}^0$ (\code{clutch\_off\_rate0}), on-rate $k_{\text{on}}$
(\code{clutch\_on\_rate}), and characteristic bond force $F_b$ (\code{bell\_force}).

\paragraph{Motor force--velocity and clutch loading.}
The actin retrograde flow slows under load through a linear force--velocity relation, and bound
clutches stretch at the difference between actin flow and substrate (fibre) motion:
\begin{equation}
  v_{\text{actin}} = v_u\,\max\!\Bigl(0,\; 1-\tfrac{T}{F_{\text{stall}}}\Bigr),
  \qquad
  \dot x_c = \max\!\bigl(0,\; v_{\text{actin}} - v_{\text{sub}}\bigr)\ \ (\text{bound}),
  \label{eq:motor}
\end{equation}
where $v_u$ = \code{unloaded\_actin\_speed}, $F_{\text{stall}}$ = the bundle's stall force
(\code{motor\_stall\_per\_protrusion}), $T=\sum_{c\in\text{bound}}k_c x_c$ is the current bundle
traction, and $v_{\text{sub}}=\mathbf v(\mathbf x)\!\cdot\!\hat{\mathbf n}_{\text{in}}$ is the
inward speed of the gripped material point (its interpolated bead velocity projected onto the
inward normal). The bundle traction $T\,\hat{\mathbf n}_{\text{in}}$ is applied to the two host
beads of the gripped edge --- the collagen is pulled \emph{inward}, toward the cell.

\paragraph{The key removal.} \GTwo\ V3 biased the two sides with a hardcoded probability: its
on-probability carried a factor $\text{side} = [\,2(1-\psi),\,2\psi\,]$ with $\psi = 0.65$
(\code{polarity\_probability}). \textbf{\GThree\ deletes this entirely}: all sectors use the same
symmetric $k_{\text{on}}$, and any binding bias is an \emph{emergent} consequence of which
protrusions physically reach and load fibres (\S\ref{sec:polarity}).

% ======================================================================
\section{What Generation-3 adds on top / G3 在其上新增的東西}

The next four sections are the actual content of \GThree. They are the \emph{only} physics that is
new relative to \S\ref{sec:engine}--\S\ref{sec:clutch}.

\subsection{A spheroid with a fibre-free gap / 帶無纖維間隙的球體}
\label{sec:gap}

\GThree\ enlarges the cell into a \emph{spheroid} of radius $R=22\um$ (vs $9\um$ in \GTwo) and leaves a
fibre-free ring of width $g$ (\code{gap}) around it. The generator rejects any candidate fibre
that intrudes inside the gap radius,
\begin{equation}
  \min_{p\in\text{fibre}}\|\mathbf p\| \;<\; R_g \equiv R + g
  \;\;\Longrightarrow\;\; \text{reject the fibre},
  \label{eq:gap}
\end{equation}
so at $t=0$ there is genuinely \textbf{no cell--ECM contact}: nothing is pre-attached, and every
fibre is a \emph{potential} contact fibre rather than a pre-placed one (contrast \GTwo, which seeds
eight curved Bézier contact fibres reaching the surface). The spheroid must \emph{earn} its grip
by reaching across the gap.

\subsection{Explicit protrusions and tip-first engagement / 顯式偽足與尖端先接觸}
\label{sec:protrusion}

The membrane is divided into $S$ equal angular sectors (\code{n\_sectors} $=24$) with outward unit
normals $\hat{\mathbf n}_k=(\cos\theta_k,\sin\theta_k)$. Sector $k$ carries one protrusion of
length $L_k$, whose \emph{tip} sits at
\begin{equation}
  \mathbf t_k = \rv_c + (R + L_k)\,\hat{\mathbf n}_k.
  \label{eq:tip}
\end{equation}

\paragraph{Length dynamics.} Every sector probes outward at a baseline speed (so the spheroid
explores in all directions and can break symmetry mechanically); sectors above the mean polarity
grow further, sectors well below it retract, and \emph{engaged} protrusions hold their length
(they are anchored). With headroom $h_k = 1 - L_k/L_{\max}$ and normalised polarity excess
$\varepsilon_k = (a_k-\bar a)/(\bar a+\epsilon)$,
\begin{equation}
  \dot L_k =
  \begin{cases}
    0, & \text{engaged},\\[3pt]
    \underbrace{v_{\text{probe}}\,h_k}_{\text{explore}}
    + \underbrace{v_{\text{grow}}\,[\varepsilon_k]_0^{3}\,h_k}_{\text{front growth}}
    - \underbrace{v_{\text{ret}}\,[-\varepsilon_k]_0^{1}}_{\text{rear retraction}}, & \text{free},
  \end{cases}
  \qquad
  L_k \leftarrow \operatorname{clip}(L_k+\dt\,\dot L_k,\,L_{\min},\,L_{\max}),
  \label{eq:length}
\end{equation}
where $[\,\cdot\,]_0^{m}=\operatorname{clip}(\cdot,0,m)$, and
$v_{\text{probe}},v_{\text{grow}},v_{\text{ret}}$ are \code{protrusion\_probe\_speed},
\code{protrusion\_growth\_speed}, \code{protrusion\_retraction\_speed}. Reach is capped at
$L_{\max}=30\um$ (\code{protrusion\_max\_length}), i.e.\ roughly the gap plus a fibre or two
($\sim\!0.7$--$4\times$ cell radius, Carey 2016; Fraley 2010/2015).

\paragraph{Tip-first engagement.} A protrusion can bind only once its tip physically reaches a
fibre. Let $d_k=\min_{\text{edges}}\operatorname{dist}(\mathbf t_k,\text{edge})$ be the distance
from the tip to the nearest material point on any fibre segment (with foot-point parameter
$\alpha_k$ and host edge $e_k$). Then
\begin{equation}
  d_k \le d_{\text{cap}}
  \;\;\Longrightarrow\;\;
  \text{sector $k$ engages edge $e_k$ at $\alpha_k$},
  \label{eq:engage}
\end{equation}
with capture distance $d_{\text{cap}}$ (\code{capture\_distance} $=2.5\um$). An engaged sector
stays engaged while any of its clutches is bound; when all release, its tip is re-tested against
the (now moved) fibre and it disengages if $d_k>d_{\text{cap}}$. Engagement is refreshed every
\code{engagement\_update\_interval} $=0.5\,\text{s}$.

\subsection{Emergent polarity: wave-pinning + adhesion feedback / 湧現極性}
\label{sec:polarity}

This is the heart of the rebuild and the direct replacement for the deleted
\code{polarity\_probability = 0.65}. The membrane carries a scalar \emph{activity}
field $a_k$ ($k=1,\dots,S$), a coarse-grained Rho-GTPase-like ``front'' marker, subject to a
\textbf{mass-conserved replicator / wave-pinning} dynamic (Mori, Jilkine \& Edelstein-Keshet 2008).
The conserved pool is
\begin{equation}
  \sum_{k=1}^{S} a_k = A_{\text{tot}} \quad(\text{\code{polarity\_total\_activity}}),
  \label{eq:mass}
\end{equation}
enforced by renormalisation each step (this is the \emph{global inhibition}). Local
self-reinforcement is a saturating Hill term plus the FAK/Rac1 adhesion signal $q_k$
(\S\ref{sec:fak}):
\begin{equation}
  \phi_k = \gamma\,\frac{a_k^{\,h}}{K^{\,h}+a_k^{\,h}} \;+\; k_{\text{FAK}}\,q_k,
  \qquad
  \text{fitness}_k = \phi_k - \langle\phi\rangle,
  \label{eq:fitness}
\end{equation}
where $\gamma$ = \code{polarity\_gamma}, $K$ = \code{polarity\_K}, $h$ = \code{polarity\_hill},
$k_{\text{FAK}}$ = \code{polarity\_fak\_gain}, and $\langle\phi\rangle$ is the sector mean. The
zero-sum fitness makes it a \textbf{replicator} (winner-take-all): sectors whose reinforcement
beats the mean grow at the expense of the rest. Membrane diffusion (a discrete ring Laplacian)
keeps the front contiguous, and a small stochastic term breaks the initial symmetry:
\begin{equation}
  (\text{lap}\,a)_k = a_{k+1} + a_{k-1} - 2a_k,
  \qquad
  a_k \;\leftarrow\; a_k
  + \frac{\dt}{\tau_p}\bigl(a_k\,\text{fitness}_k + D\,(\text{lap}\,a)_k\bigr)
  + \sigma\sqrt{\dt}\,\xi_k,
  \label{eq:polarity}
\end{equation}
followed by $a_k\leftarrow\max(a_k,10^{-4})$ and renormalisation to
Eq.~\eqref{eq:mass}. Here $\tau_p$ = \code{polarity\_time} (front relaxation time),
$D$ = \code{polarity\_diffusion}, $\sigma$ = \code{polarity\_noise}, and
$\xi_k\sim\mathcal N(0,1)$ i.i.d.\ per step.

\paragraph{How it works.} From a near-uniform, symmetric start
($a_k \approx A_{\text{tot}}/S$ with $2\%$ noise), the bistable local kinetics plus mass
conservation collapse the field to \textbf{one stable broad front} (a pinned wave), rather than a
spread of many peaks. Three ingredients play distinct roles:
\begin{itemize}[leftmargin=1.4em,itemsep=2pt]
  \item \textbf{Noise $\sigma\sqrt{\dt}\,\xi_k$} chooses \emph{which} direction --- the sign is
        genuinely stochastic and differs between seeds (verified: five seeds polarise toward
        $84^\circ,\,-84^\circ,\,-8^\circ,\,-99^\circ,\,64^\circ$).
  \item \textbf{Adhesion feedback $k_{\text{FAK}}q_k$} pulls the front toward sectors that
        \emph{actually grip} collagen --- a mechanochemical positive loop.
  \item \textbf{Aligned collagen} (optional fixture) biases only the front \emph{axis}, never its
        sign.
\end{itemize}
Nothing prescribes $+x$, a left/right probability, or the number $0.65$.

\subsection{FAK/Rac1 adhesion signal / 黏著訊號}
\label{sec:fak}

The feedback variable $q_k$ is a low-pass-filtered measure of adhesion success in sector $k$. Let
$n_k$ be the number of currently bound clutches and $T_k$ the bundle traction. Define an
engagement saturation and a normalised load, combine them, and filter:
\begin{equation}
  q_k^{\text{raw}} = \underbrace{\bigl(1-e^{-n_k/n_s}\bigr)}_{\text{engagement}}\cdot
  \underbrace{\frac{\lambda_k}{1+\lambda_k}}_{\text{load}},
  \qquad
  \lambda_k = \frac{T_k}{\tfrac12 F_{\text{stall}}},
  \qquad
  q_k \leftarrow q_k + \frac{\dt}{\tau_q}\bigl(q_k^{\text{raw}}-q_k\bigr),
  \label{eq:fak}
\end{equation}
with clutch scale $n_s$ = \code{adhesion\_clutch\_scale} and filter time $\tau_q$ =
\code{adhesion\_filter\_time}. The low-pass ($\tau_q=20\,\text{s}$) makes the feedback respond to
\emph{sustained} traction rather than momentary binding flicker, which is what stabilises a single
front instead of the ``twitchy'' flicker of the superseded \GThree.

\subsection{Reaction-driven 2-D migration / 反作用力驅動的二維遷移}
\label{sec:motion}

Motion follows the same reaction principle as \GTwo\ V3 but in full 2-D. The clutch bundles apply a
total inward force field $\Fv^{\text{active}}$ on the collagen (\S\ref{sec:clutch}). Because those
clutch forces are internal to the cell--ECM pair, the net reaction on the rigid spheroid is
minus their sum, and it points \emph{toward} the gripping front:
\begin{equation}
  \mathbf R = -\sum_i \Fv_i^{\text{active}},
  \qquad
  \mathbf v_c = \frac{\mathbf R}{\gamma_c},
  \qquad
  \mathbf v_c \leftarrow \mathbf v_c\,\frac{\min(\|\mathbf v_c\|,\,v_{\max})}{\|\mathbf v_c\|},
  \qquad
  \rv_c \leftarrow \rv_c + \dt\,\mathbf v_c,
  \label{eq:motion}
\end{equation}
with cell drag $\gamma_c$ = \code{cell\_drag} $=3000\,\text{nN\,s/}\umt$ (calibrated to
$0.1$--$0.3\um$/min) and a numerical/biological speed guard $v_{\max}$ = \code{max\_cell\_speed}.
In \code{--fixed} mode the spheroid is clamped ($\mathbf v_c\equiv0$): fibres still get pulled and
realigned, but there is no migration --- a clean falsifiability control.

\emph{Contrast with \GTwo\ V3:} V3 intentionally allowed only one translational degree of freedom,
$\dot x_c = \text{reaction}_x/\gamma_c$, $\dot y_c = 0$, on a declared left/right axis. \GThree\ has no
declared axis: both components of $\mathbf R$ act, and the direction is set by the emergent front.

% ======================================================================
\section{The full G3 time step / G3 完整時間步}
\label{sec:loop}

Each step of \code{run\_spheroid} (size $\dt$) executes the following, in order. Steps 1--3 and 5
are new; step 4 is the shared clutch law (\S\ref{sec:clutch}); step 7 is the shared engine
(\S\ref{sec:engine}).

\begin{enumerate}[leftmargin=1.6em,itemsep=2pt]
  \item \textbf{Emergent polarity} --- one wave-pinning update, Eq.~\eqref{eq:polarity}
        (\code{step\_polarity}).
  \item \textbf{Protrusion growth} --- update every $L_k$, Eq.~\eqref{eq:length}
        (\code{step\_protrusions}).
  \item \textbf{Tip-first engagement} --- every $0.5\,\text{s}$, re-test tips and (dis)engage,
        Eq.~\eqref{eq:engage} (\code{update\_engagement}).
  \item \textbf{Motor--clutch loading} --- advance every bundle, build the inward nodal force
        field $\Fv^{\text{active}}$, Eqs.~\eqref{eq:bell}--\eqref{eq:motor}
        (\code{step\_clutches}).
  \item \textbf{Adhesion filter} --- update $q_k$, Eq.~\eqref{eq:fak}
        (\code{update\_adhesion}).
  \item \textbf{Reaction-driven motion} --- move the spheroid, Eq.~\eqref{eq:motion}.
  \item \textbf{Advance the collagen} --- one overdamped step of the frozen \GTwo\ engine,
        Eqs.~\eqref{eq:balance}--\eqref{eq:repulsion} (\code{fast\_advance}); returns the bead
        velocity field used as $v_{\text{sub}}$ next step.
\end{enumerate}

The narrative this produces: $t=0$ fibre-free gap $\to$ protrusions probe outward in all
directions $\to$ some tips reach fibres and bind $\to$ motor--clutches load and pull those fibres
inward (the network realigns radially at the front) $\to$ a single front is selected and reinforced
by adhesion feedback $\to$ the spheroid migrates toward that front, remodelling collagen along its
path.

% ======================================================================
\section{Metrics / 量測}
\label{sec:metrics}

\paragraph{Radial nematic order.} For each fibre segment with unit tangent $\hat{\mathbf t}$ and
midpoint radial unit vector $\hat{\mathbf e}_r$ (from the spheroid centre), the nematic radial order
is
\begin{equation}
  S_r = \bigl\langle\,2(\hat{\mathbf t}\!\cdot\!\hat{\mathbf e}_r)^2 - 1\,\bigr\rangle,
\end{equation}
averaged in near/mid/far shells (surface distances $0$--$20$, $20$--$45$, $45$--$90\um$).
$S_r=+1$ means fibres point radially (toward the spheroid), $0$ random, $-1$ tangential. Radial
pulling drives the near-shell $S_r$ from $\approx\!-0.62$ toward $-0.31$ over the run.

\paragraph{Polarity vector and traction.} The net polarity direction is
$\mathbf P=\sum_k a_k\hat{\mathbf n}_k$ (its angle is the migration heading); the collective
traction is $\sum_k T_k$; and the path length accumulates $\|\rv_c(t+\dt)-\rv_c(t)\|$.

% ======================================================================
\section{Parameters and literature anchors / 參數與文獻依據}
\label{sec:params}

\begin{longtable}{@{}l l r l@{}}
\toprule
\textbf{Symbol} & \textbf{Config field} & \textbf{Value} & \textbf{Meaning / source} \\
\midrule
\endhead
\multicolumn{4}{@{}l}{\emph{Shared \GTwo\ engine}}\\
$L$ & \code{domain\_size} & $240\um$ & field size (\GTwo: $180$) \\
$E$ & \code{collagen\_modulus\_mpa} & $32\,\text{MPa}$ & wet fibre modulus (Lee 2014) \\
$D$ & \code{fiber\_diameter} & $0.30\um$ & fibre diameter \\
$\rho$ & \code{compression\_ratio} & $0.10$ & microbuckling softening (Notbohm 2015) \\
$k_x$ & \code{crosslink\_stiffness} & $75\,\text{nN}/\umt$ & hinged-link penalty \\
$k_{\text{rep}}$ & \code{repulsion\_stiffness} & $600\,\text{nN}/\umt$ & cell-surface penalty \\
$\zeta$ & \code{bead\_drag} & $360\,\text{nN\,s}/\umt$ & bead drag (\GTwo\ V3) \\
$\dt$ & \code{dt} & $0.05\,\text{s}$ & integration step \\
\midrule
\multicolumn{4}{@{}l}{\emph{Spheroid \& protrusions (new)}}\\
$R$ & \code{cell\_radius} & $22\um$ & spheroid radius (\GTwo: $9$) \\
$g$ & \code{gap} & $12\um$ & fibre-free ring \\
$N_f$ & \code{n\_fibers} & $150$ & fibre count (\GTwo: $99$) \\
$S$ & \code{n\_sectors} & $24$ & membrane sectors / protrusions \\
$L_{\max}$ & \code{protrusion\_max\_length} & $30\um$ & max reach (Carey 2016) \\
$v_{\text{probe}}$ & \code{protrusion\_probe\_speed} & $0.14\um/\text{s}$ & baseline probing \\
$v_{\text{grow}}$ & \code{protrusion\_growth\_speed} & $0.16\um/\text{s}$ & front growth \\
$v_{\text{ret}}$ & \code{protrusion\_retraction\_speed} & $0.22\um/\text{s}$ & rear retraction \\
$d_{\text{cap}}$ & \code{capture\_distance} & $2.5\um$ & tip capture radius \\
\midrule
\multicolumn{4}{@{}l}{\emph{Motor--clutch bundle (shared law)}}\\
$C$ & \code{n\_clutches\_per\_protrusion} & $8$ & clutches per bundle \\
$k_c$ & \code{clutch\_stiffness} & $2.0\,\text{nN}/\umt$ & clutch spring \\
$k_{\text{on}}$ & \code{clutch\_on\_rate} & $0.08\,\text{s}^{-1}$ & symmetric binding \\
$k_{\text{off}}^0$ & \code{clutch\_off\_rate0} & $0.02\,\text{s}^{-1}$ & unloaded off-rate \\
$F_b$ & \code{bell\_force} & $2.0\nN$ & Bell bond force \\
$v_u$ & \code{unloaded\_actin\_speed} & $0.03\um/\text{s}$ & unloaded actin flow \\
$F_{\text{stall}}$ & \code{motor\_stall\_per\_protrusion} & $4.0\nN$ & motor stall \\
\midrule
\multicolumn{4}{@{}l}{\emph{Emergent polarity + adhesion (new)}}\\
$A_{\text{tot}}$ & \code{polarity\_total\_activity} & $6.0$ & conserved pool \\
$\gamma$ & \code{polarity\_gamma} & $2.2$ & autocatalytic gain \\
$K$ & \code{polarity\_K} & $0.5$ & Hill half-saturation \\
$h$ & \code{polarity\_hill} & $2.0$ & Hill cooperativity \\
$D$ & \code{polarity\_diffusion} & $0.45$ & membrane smoothing \\
$\sigma$ & \code{polarity\_noise} & $0.08$ & symmetry-breaking noise \\
$k_{\text{FAK}}$ & \code{polarity\_fak\_gain} & $2.5$ & FAK/Rac1 feedback gain \\
$\tau_p$ & \code{polarity\_time} & $6.0\,\text{s}$ & front relaxation time \\
$\tau_q$ & \code{adhesion\_filter\_time} & $20\,\text{s}$ & $q_k$ low-pass \\
$n_s$ & \code{adhesion\_clutch\_scale} & $3.0$ & engagement scale \\
\midrule
\multicolumn{4}{@{}l}{\emph{Motion (new geometry, shared principle)}}\\
$\gamma_c$ & \code{cell\_drag} & $3000\,\text{nN\,s}/\umt$ & spheroid drag \\
$v_{\max}$ & \code{max\_cell\_speed} & $0.02\um/\text{s}$ & speed guard \\
\bottomrule
\end{longtable}

\noindent The literature sets \emph{targets and plausible scales}, not fits: spheroid traction of
tens of nN (Steinwachs 2016; Mark 2020), MDA-MB-231 3-D speed $0.1$--$0.3\um$/min (Sapudom 2019;
Aguilar-Rojas 2024), protrusion reach $\sim\!0.7$--$4\times$ cell radius (Carey 2016; Fraley
2010/2015), wave-pinning polarisation in $\sim\!10\,\text{s}$ with stochastic sign (Mori 2008;
Jilkine 2011), and fibre realignment over tens of minutes to hours (Kim 2017; Han 2018).

% ======================================================================
\section{Side-by-side: G2 vs G3 / 逐項對照}
\label{sec:compare}

\renewcommand{\arraystretch}{1.25}
\begin{longtable}{@{}p{0.24\textwidth} p{0.35\textwidth} p{0.35\textwidth}@{}}
\toprule
\textbf{Aspect} & \textbf{\GTwo\ (corrected)} & \textbf{\GThree\ (spheroid guidance)} \\
\midrule
\endhead
Collagen engine & bead--spring, overdamped, axial $EA/\ell_0$, bending $EI/\ell^3$, hinged crosslinks, boundary anchoring, soft repulsion & \textbf{identical} --- reused verbatim (\S\ref{sec:engine}) \\
Integrator & \code{Network.advance} (\code{np.add.at}) & \code{fast\_advance} (\code{np.bincount}); \textbf{bit-identical}, $\sim2.4\times$ faster \\
Cell body & rigid circle, $R=9\um$ & rigid spheroid, $R=22\um$ \\
Contact at $t=0$ & 8 pre-placed Bézier contact fibres touch the surface & \textbf{fibre-free gap} --- no contact; nothing pre-attached (Eq.~\ref{eq:gap}) \\
Cell--ECM coupling & Gaussian-weighted contact patches in a $\pm30^\circ$ sector & \textbf{explicit protrusions}, tip-first engagement (Eqs.~\ref{eq:tip}--\ref{eq:engage}) \\
Motor--clutch law & Bell slip-bond $+$ motor F--V (V3) & \textbf{same law}, one bundle per protrusion (\S\ref{sec:clutch}) \\
Directional bias & hardcoded \code{polarity\_probability}$=0.65$ (V3 side factor) & \textbf{removed}; emergent wave-pinning $+$ FAK feedback (\S\ref{sec:polarity}) \\
Polarity origin & prescribed constant & noise breaks symmetry; mechanics pick direction; sign is stochastic per seed \\
Migration & 1-D, $\dot x_c=\text{reaction}_x/\gamma_c$, $\dot y_c=0$ (V3) & \textbf{full 2-D}, $\mathbf v_c=\mathbf R/\gamma_c$ (Eq.~\ref{eq:motion}) \\
Force scale & $5\nN$ local pull (V2) & tens of nN collective ($\sim\!36\nN$ peak) \\
Domain / fibres & $180\um$ / 99 fibres & $240\um$ / 150 fibres \\
Percolation gate & connected fraction $\ge0.85$ & $\ge0.80$ (relaxed for the gap) \\
Duration & $30$--$60\,\text{min}$ & $60\,\text{min}$ (realign $+$ migrate) \\
Constitutive law & purely elastic; no SLS, no plasticity & \textbf{identical} --- no SLS, no plasticity \\
\bottomrule
\end{longtable}

\noindent\textbf{One-line summary.} \GThree\ = \emph{the \GTwo\ collagen engine and \GTwo\ motor--clutch
law, unchanged}, wrapped around a spheroid-with-gap geometry, explicit tip-first protrusions, an
emergent wave-pinning polarity field that \emph{replaces} the $0.65$ constant, and full-2-D
reaction-driven motion.

% ======================================================================
\section{Reference run and falsifiability / 參考結果與可證偽性}
\label{sec:results}

The reference run --- 150 fibres, $240\um$ field, $22\um$ spheroid, $12\um$ gap, $60\,\text{min}$,
seed 7 (\code{results/g3\_spheroid\_guidance/main/}) --- gives:

\begin{center}
\begin{tabular}{@{}l r l@{}}
\toprule
\textbf{Quantity} & \textbf{Value} & \textbf{Literature target} \\
\midrule
Peak collective traction & $\approx 36\nN$ & tens of nN, $\le100$ (Steinwachs 2016) \\
Net migration in $60\,$min & $\approx 18.8\um$ & --- \\
Migration speed & $\approx 0.32\um$/min & $0.1$--$0.3\um$/min (Sapudom 2019) \\
Engaged protrusions at end & $\approx 10$ & --- \\
Near-shell radial order & $-0.62 \to -0.31$ & realignment over tens of min--h \\
Max bead displacement & $\approx 5.1\um$ & --- \\
Connected fraction & $1.0$ & gate $\ge 0.80$ \\
\bottomrule
\end{tabular}
\end{center}

\noindent\textbf{Emergence check.} Across five seeds the migration heading is
$84^\circ,\,-84^\circ,\,-8^\circ,\,-99^\circ,\,64^\circ$ (spread $\approx75^\circ$): the direction
is genuinely stochastic and set per realisation, confirming there is no built-in bias.

\noindent\textbf{Controls that make the mechanism falsifiable.}
\begin{itemize}[leftmargin=1.4em,itemsep=1pt]
  \item \code{--fixed} clamps the spheroid: fibres still realign, but migration vanishes.
  \item Zero motor stall force: protrusions attach without deforming fibres.
  \item An empty gap with no reachable fibre: zero engagement, zero traction, zero motion.
\end{itemize}

Two view modes are rendered in the \GTwo\ visual grammar (springs coloured by strain, gold-diamond
crosslinks, black-square fixed beads, peach cell, orange protrusions/contacts, faint ghost of the
initial network, true $1\times$ scale $+$ scale bar): \code{*\_full.gif} (full field) and
\code{*\_follow.gif} (follow-cell zoom), matching the G2-v3 notebook toggle.

% ======================================================================
\section{Evidence boundary / 證據邊界}

\GThree\ is a \textbf{2-D minimal mechanism model of emergent cell--collagen guidance}. It is
\emph{not} a realistic three-dimensional tumour-migration prediction: it contains no calibrated
collagen concentration or pore topology, no deformable cell or nucleus, no proteolysis, no EMT, no
cell--cell interactions, no stress relaxation, and no irreversible plasticity. \GTwo\
(\code{generations/g2\_corrected/}) remains frozen and unchanged; \GThree\ reuses its engine without
rewriting any \GTwo\ file. Every generated figure carries the boundary:
\emph{``Mechanism demo --- not a 3-D tumour-migration prediction.''}

% ======================================================================
\section*{References / 參考文獻}
\addcontentsline{toc}{section}{References}
\small
\begin{enumerate}[leftmargin=1.6em,itemsep=1pt]
  \item Lee et al.\ (2014). A 3-D computational model of collagen network mechanics. \emph{PLoS ONE} 9:e111896.
  \item Notbohm et al.\ (2015). Microbuckling of fibrin provides a mechanism for cell mechanosensing. \emph{J. R. Soc. Interface}.
  \item Bell (1978). Models for the specific adhesion of cells to cells. \emph{Science} 200:618.
  \item Bangasser \& Odde (2013). Master-equation analysis of a motor-clutch model. \emph{Cell. Mol. Bioeng.} 6:449.
  \item Prahl et al.\ (2020). Predicting confined 1-D cell migration from a 2-D motor-clutch model. \emph{Biophys. J.} 118:1709.
  \item Steinwachs et al.\ (2016). Three-dimensional force microscopy of cells in biopolymer networks. \emph{Nat. Methods} 13:171.
  \item Carey et al.\ (2016). Local ECM alignment directs protrusion dynamics and migration through Rac1 and FAK. \emph{Integr. Biol.} 8:821.
  \item Mori, Jilkine \& Edelstein-Keshet (2008). Wave-pinning and cell polarity from a bistable reaction--diffusion system. \emph{Biophys. J.} 94:3684.
  \item Fraley et al.\ (2010). A distinctive role for focal-adhesion proteins in 3-D motility. \emph{Nat. Cell Biol.} 12:598.
  \item Fraley et al.\ (2015). 3-D matrix fibre alignment modulates migration and MT1-MMP utility. \emph{Sci. Rep.} 5:14580.
  \item Kim et al.\ (2017). Stress-induced plasticity of dynamic collagen networks. \emph{Nat. Commun.} 8:842.
  \item Han et al.\ (2018). Cell contraction induces long-ranged stress stiffening in the ECM. \emph{PNAS}.
  \item Sapudom et al.\ (2019) / Aguilar-Rojas et al.\ (2024). MDA-MB-231 3-D migration speeds in collagen.
\end{enumerate}
\normalsize

\vfill
\noindent\footnotesize\emph{Source of truth:} \code{generations/g3\_spheroid\_guidance/model.py}
(\code{SpheroidConfig}, \code{step\_polarity}, \code{step\_protrusions}, \code{update\_engagement},
\code{step\_clutches}, \code{update\_adhesion}, \code{run\_spheroid}, \code{fast\_advance}) and the
frozen \GTwo\ engine \code{generations/g2\_corrected/common/model.py}. All equations above were
transcribed directly from that code.

\end{document}
